\chapter{Dalsze twierdzenia graniczne i koncentracja prawdopodobieństwa}
\label{ch:advancedlimits}

\section{Od praw asymptotycznych do skończonych prób}

Prawo wielkich liczb odpowiada na pytanie, czy średnia
stabilizuje się w granicy $n\to\infty$, a CTG opisuje
kształt rozkładu przeskalowanych błędów. W praktyce
często chcemy wiedzieć więcej: jak szybko zanika
prawdopodobieństwo dużego odchylenia, co zrobić, jeśli
składniki mają różne rozkłady, albo czy zbieżność
średniej pozostaje prawdziwa po nieliniowym
przekształceniu. Wprowadzimy narzędzia
odpowiadające na te pytania, zachowując pełne
dowody w przedstawianych wersjach twierdzeń.

\section{Lemat Hoeffdinga i koncentracja średniej}

\begin{lemat}[Nierówność na wariancję zmiennej ograniczonej]
Niech $Z$ będzie rzeczywistą zmienną losową taką,
że $a\leq Z\leq b$ p.n., gdzie $a<b$.
Wówczas $\Var(Z)\leq(b-a)^2/4$.
\end{lemat}
\begin{proof}
Ponieważ $(Z-a)(b-Z)\geq0$, po rozwinięciu
i uśrednieniu otrzymujemy
$\E Z^2\leq(a+b)\E Z-ab$.
Niech $m=\E Z$. Wobec tego
$\Var(Z)=\E Z^2-m^2\leq(b-m)(m-a)$.
Iloczyn dwóch nieujemnych liczb o sumie $b-a$
nie przekracza $(b-a)^2/4$. To dowodzi tezy.
\end{proof}

\begin{lemat}[Lemat Hoeffdinga]
Niech $Z$ będzie rzeczywistą zmienną losową,
$a\leq Z\leq b$ p.n. dla pewnych $a<b$.
Dla każdego $t\in\R$ zachodzi
\[
\E\exp(t(Z-\E Z))
\leq\exp\left(\frac{t^2(b-a)^2}{8}\right).
\]
\end{lemat}
\begin{proof}
Zdefiniujmy $M(t)=\E e^{tZ}$ i
$h(t)=\log M(t)$. Ponieważ $Z$ jest ograniczona, $M(t)$ jest skończona
i ściśle dodatnia dla każdego $t$. Jeśli
$B=\max(|a|,|b|)$, to dla $|s-t|\leq1$ mamy
$|Z^j e^{sZ}|\leq B^j e^{(|t|+1)B}$ p.n.
Z twierdzenia o różniczkowaniu pod znakiem całki
(lub z twierdzenia o zbieżności zmajoryzowanej
dla ilorazów różnicowych) otrzymujemy
$M'(t)=\E[Ze^{tZ}]$ oraz $M''(t)=\E[Z^2e^{tZ}]$. Obliczamy
$h'(t)=M'(t)/M(t)$ oraz
$h''(t)=M''(t)/M(t)-(M'(t)/M(t))^2$.
Dla każdego $t$ wyrażenie $h''(t)$ jest
wariancją $Z$ względem nowej miary (przechylenie wykładnicze)
$\dd\Prob_t=e^{tZ}\dd\Prob/M(t)$. Nowa miara jest
probabilistyczna, ponieważ jej gęstość jest nieujemna
i całkuje się do jedności. Ponadto $\Prob_t\ll\Prob$,
więc $a\leq Z\leq b$ również $\Prob_t$-prawie na pewno.
Rzeczywiście $M'(t)/M(t)=\E_t Z$ oraz
$M''(t)/M(t)=\E_t Z^2$, gdzie $\E_t$ oznacza
wartość oczekiwaną względem $\Prob_t$.
Poprzedni lemat daje więc
$0\leq h''(t)\leq(b-a)^2/4$ dla wszystkich $t$.

Funkcja $g(t)=h(t)-th'(0)$ spełnia $g(0)=g'(0)=0$,
a $g''(t)=h''(t)\leq(b-a)^2/4$.
Stosując wzór Taylora z resztą całkową do funkcji
$s\mapsto g(st)$ na $[0,1]$, otrzymujemy dla
dowolnego $t\in\R$
\[
g(t)=t^2\int_0^1(1-s)g''(st)\,\dd\lambda_1(s)
\leq \frac{t^2(b-a)^2}{8}.
\]
Wzór obowiązuje także dla $t<0$, ponieważ czynnik
$t^2(1-s)$ jest nieujemny.
Ponieważ $h'(0)=\E Z$,
$g(t)=\log\E e^{t(Z-\E Z)}$.
Potęgowanie daje tezę.
\end{proof}

\begin{twierdzenie}[Nierówność Hoeffdinga dla niezależnych składników]
Niech $X_1,\ldots,X_n$ będą niezależnymi rzeczywistymi
zmiennymi losowymi oraz niech $a_k\leq X_k\leq b_k$
p.n., gdzie $a_k<b_k$. Oznaczmy
$S_n=\sum_{k=1}^n(X_k-\E X_k)$.
Dla każdego $u>0$ zachodzą nierówności
\[
\Prob(S_n\geq u)\leq
\exp\left(-\frac{2u^2}{\sum_{k=1}^n(b_k-a_k)^2}\right),
\qquad
\Prob(|S_n|\geq u)\leq
2\exp\left(-\frac{2u^2}{\sum_{k=1}^n(b_k-a_k)^2}\right).
\]
\end{twierdzenie}
\begin{proof}
Dla $t>0$ nierówność Markowa zastosowana
do $e^{tS_n}$ daje
$\Prob(S_n\geq u)\leq e^{-tu}\E e^{tS_n}$.
Z niezależności i lematu Hoeffdinga
$\E e^{tS_n}=\prod_k\E e^{t(X_k-\E X_k)}
\leq\exp(t^2V/8)$,
gdzie $V=\sum_k(b_k-a_k)^2$.
Otrzymujemy oszacowanie
$\Prob(S_n\geq u)\leq\exp(-tu+t^2V/8)$.
Minimalizując kwadratową funkcję wykładnika po $t>0$,
wybieramy $t=4u/V$, skąd $-tu+t^2V/8=-2u^2/V$.
Tę samą nierówność stosujemy do $-S_n$ i korzystamy
z sumowania prawdopodobieństw dwóch ogonów.
\end{proof}

\begin{przyklad}[Ile prób potrzeba do kontroli częstości?]
Jeżeli $B_1,\ldots,B_n$ są niezależnymi próbami
Bernoulliego z parametrem $p\in[0,1]$, to $0\leq B_k\leq1$.
Dla $\overline B_n=n^{-1}\sum_{k=1}^n B_k$ mamy
$\E\overline B_n=p$ z liniowości wartości oczekiwanej. Nierówność Hoeffdinga
przy $u=n\varepsilon$ daje
$\Prob(|\overline B_n-p|\geq\varepsilon)
\leq2e^{-2n\varepsilon^2}$.
Dla $\varepsilon>0$ i zadanej $\delta\in(0,1)$,
aby prawa strona była nie większa niż $\delta$, wystarczy
$n\geq\log(2/\delta)/(2\varepsilon^2)$.
Na przykład dla $\varepsilon=0{,}05$ i $\delta=0{,}05$
wystarczy $n=\lceil 200\log 40\rceil=738$ prób,
niezależnie od nieznanego parametru $p$.
Jest to gwarancja nieasymptotyczna, ważna dla
każdego $n$, choć nie zawsze ostra.
\end{przyklad}

Zależność ograniczenia od liczby prób i progu
odchylenia ilustruje rysunek~\ref{fig:rp1-11-hoeffding}.

\begin{figure}[htbp]
\centering
\begin{tikzpicture}[x=22cm,y=5cm,>=Latex]
\draw[->] (0,0)--(.44,0) node[right] {$\varepsilon$};
\draw[->] (0,0)--(0,1.08) node[above] {ograniczenie};
\foreach \x/\lab in {.1/{0{,}1},.2/{0{,}2},.3/{0{,}3},.4/{0{,}4}}
  {\draw (\x,0)--++(0,-.018) node[below] {$\lab$};}
\foreach \y/\lab in {.25/{0{,}25},.5/{0{,}5},.75/{0{,}75},1/{1}}
  {\draw[gray!20] (0,\y)--(.4,\y);
   \draw (0,\y)--++(-.005,0) node[left] {$\lab$};}
\draw[blue!75!black,very thick,domain=0:.4,samples=121]
  plot (\x,{min(1,2*exp(-40*\x*\x))});
\draw[teal!75!black,very thick,dashed,domain=0:.4,samples=121]
  plot (\x,{min(1,2*exp(-100*\x*\x))});
\draw[orange!85!black,very thick,dashdotted,domain=0:.4,samples=121]
  plot (\x,{min(1,2*exp(-200*\x*\x))});
\draw[blue!75!black,very thick] (.29,.91)--(.325,.91)
  node[right,black] {$n=20$};
\draw[teal!75!black,very thick,dashed] (.29,.80)--(.325,.80)
  node[right,black] {$n=50$};
\draw[orange!85!black,very thick,dashdotted] (.29,.69)--(.325,.69)
  node[right,black] {$n=100$};
\end{tikzpicture}
\caption{Dwustronne ograniczenie Hoeffdinga
$\min(1,2e^{-2n\varepsilon^2})$ dla średniej prób
Bernoulliego. Krzywe są ograniczeniami górnymi
prawdopodobieństwa odchylenia, a nie dokładnymi
prawdopodobieństwami. Przy ustalonym $\varepsilon>0$
ograniczenie jest nierosnące względem $n$.}
\label{fig:rp1-11-hoeffding}
\end{figure}

\begin{twierdzenie}[Prawie pewna szybkość stabilizacji średnich]
Niech $X_1,X_2,\ldots$ będą niezależnymi rzeczywistymi
zmiennymi losowymi o wspólnej średniej $m$ oraz
$a\leq X_k\leq b$ p.n. dla wszystkich $k$,
gdzie $a<b$. Oznaczmy $\overline X_n=n^{-1}\sum_{k=1}^n X_k$.
Dla każdego $\eta>0$ zachodzi
\[
|\overline X_n-m|\leq
(b-a)\sqrt{\frac{(1+\eta)\log n}{2n}}
\quad\text{dla wszystkich dostatecznie dużych }n
\text{ prawie na pewno}.
\]
\end{twierdzenie}
\begin{proof}
Wstawiamy
$\varepsilon_n=(b-a)\sqrt{(1+\eta)\log n/(2n)}$
do dwustronnego oszacowania Hoeffdinga. Dla $n\geq2$
\[
\Prob(|\overline X_n-m|\geq\varepsilon_n)
\leq2e^{-2n\varepsilon_n^2/(b-a)^2}
=2n^{-(1+\eta)}.
\]
Szereg tych prawdopodobieństw jest zbieżny.
Pierwszy lemat Borela--Cantellego
(rozdział~\ref{ch:probability}) stwierdza, że
zdarzenia $\{|\overline X_n-m|\geq\varepsilon_n\}$
zachodzą tylko skończenie wiele razy p.n.
Na zbiorze pełnej miary od pewnego (zależnego
od wyniku doświadczenia) indeksu mamy więc
$|\overline X_n-m|<\varepsilon_n$.
\end{proof}

\begin{uwaga}
Oszacowanie rzędu $\sqrt{(\log n)/n}$ jest
wystarczającym, prostym w dowodzie tempem
prawie pewnej zbieżności. Nie należy utożsamiać
go z optymalnymi, bardziej subtelnymi wynikami
o fluktuacjach pojedynczych trajektorii.
\end{uwaga}

\section{Centralne twierdzenie graniczne dla niejednakowych składników}

Warunek jednakowego rozkładu w podstawowym CTG
można osłabić. Problemem stają się wówczas duże,
rzadkie składniki: jeden z nich może zdominować
całą sumę i zniszczyć granicę normalną.
Warunek Lindeberga wymusza, aby ich wkład do
wariancji po odpowiednim skalowaniu zanikał.

\begin{definicja}[Schemat trójkątny i warunek Lindeberga]
Dla $n\geq1$ niech $k_n\in\N$, a
$Y_{n,1},\ldots,Y_{n,k_n}$ niech będą
niezależnymi \emph{w obrębie każdego wiersza}
rzeczywistymi zmiennymi losowymi spełniającymi
$\E Y_{n,k}=0$ i $\E Y_{n,k}^2<\infty$.
Nie zakładamy niezależności między różnymi wierszami.
Połóżmy
$s_n^2=\sum_{k=1}^{k_n}\E Y_{n,k}^2>0$.
Mówimy, że zachodzi \emph{warunek Lindeberga}, jeżeli
 dla każdego $\varepsilon>0$
\[
\frac1{s_n^2}\sum_{k=1}^{k_n}
\E\left[Y_{n,k}^2\chi_{\{|Y_{n,k}|>\varepsilon s_n\}}\right]
\longrightarrow0.
\]
\end{definicja}

\begin{twierdzenie}[CTG w wersji Lindeberga: kierunek wystarczający]
Jeżeli schemat trójkątny spełnia powyższe założenia
oraz warunek Lindeberga, to
\[
\frac{1}{s_n}\sum_{k=1}^{k_n}Y_{n,k}
\xrightarrow{d}\Normal(0,1).
\]
\end{twierdzenie}
\begin{proof}
\emph{Krok 1: normalizacja i znikanie największej
wariancji.} Połóżmy $Z_{n,k}=Y_{n,k}/s_n$.
Wtedy $\sum_k\E Z_{n,k}^2=1$, a warunek Lindeberga
przyjmuje postać
$\sum_k\E[Z_{n,k}^2\chi_{\{|Z_{n,k}|>\varepsilon\}}]\to0$.
Dla dowolnego $\varepsilon>0$
\[
\E Z_{n,k}^2\leq\varepsilon^2+
\E[Z_{n,k}^2\chi_{\{|Z_{n,k}|>\varepsilon\}}],
\]
zatem $\limsup_n\max_k\E Z_{n,k}^2\leq\varepsilon^2$.
Ponieważ $\varepsilon$ jest dowolne,
$\max_k\E Z_{n,k}^2\to0$. To jest warunek,
że żadna pojedyncza zmienna nie niesie zauważalnej
części całkowitej wariancji.

\emph{Krok 2: rozwinięcie pojedynczych funkcji
charakterystycznych.} Ustalmy $t\in\R$.
Dla każdej liczby rzeczywistej $u$ wzór Taylora
z resztą całkową, zastosowany do funkcji
$s\mapsto e^{isu}$ na $[0,1]$, daje
\[
e^{iu}-1-iu=-u^2\int_0^1(1-s)e^{isu}\,\dd\lambda_1(s),
\]
więc $|e^{iu}-1-iu|\leq u^2/2$. Odejmując
odpowiedni wyraz drugiego rzędu i korzystając
z $|e^{isu}-1|\leq s|u|$, otrzymujemy
$|e^{iu}-1-iu+u^2/2|\leq |u|^3/6$.
Użyta nierówność wynika z
$e^{isu}-1=iu\int_0^s e^{ivu}\,\dd\lambda_1(v)$.
Z pierwszego oszacowania i nierówności trójkąta
ten sam moduł jest nie większy niż $u^2$.
Łącznie
$|e^{iu}-1-iu+u^2/2|\leq u^2\min(|u|,1)$.
Oznaczmy $\psi_{n,k}(t)=\E e^{itZ_{n,k}}$.
Ponieważ $\E Z_{n,k}=0$, dostajemy
\[
\psi_{n,k}(t)=1-\frac{t^2}{2}\E Z_{n,k}^2+r_{n,k}(t),
\qquad
|\psi_{n,k}(t)-1|\leq\frac{t^2}{2}\E Z_{n,k}^2.
\]
Pozostaje sprawdzić, że $\sum_k|r_{n,k}(t)|\to0$.
Dla ustalonego $\delta>0$ rozbijmy oczekiwanie
w reszcie na $\{|Z_{n,k}|\leq\delta\}$ i dopełnienie.
Na pierwszej części, ponieważ
$\min(|tZ_{n,k}|,1)\leq |t|\delta$,
otrzymujemy oszacowanie $|t|^3\delta\E Z_{n,k}^2$.
Na drugiej części korzystamy z
$\min(|tZ_{n,k}|,1)\leq1$, uzyskując
$t^2\E[Z_{n,k}^2\chi_{\{|Z_{n,k}|>\delta\}}]$.
Możemy więc przyjąć $C_t=\max(|t|^3,t^2)$. Po zsumowaniu
i wykorzystaniu sumy wariancji równej jeden,
\[
\limsup_n\sum_k|r_{n,k}(t)|\leq C_t\delta.
\]
Przechodzimy do granicy przy $\delta\downarrow0$, uzyskując
$\sum_k|r_{n,k}(t)|\to0$, a więc
$\sum_k(\psi_{n,k}(t)-1)\to-t^2/2$.

\emph{Krok 3: od sumy do iloczynu.} Niech
$z_{n,k}=\psi_{n,k}(t)-1$. Z poprzednich oszacowań
$\max_k|z_{n,k}|\to0$ oraz
$\sum_k|z_{n,k}|\leq t^2/2$.
W konsekwencji
$\sum_k|z_{n,k}|^2\leq
(\max_k|z_{n,k}|)\sum_k|z_{n,k}|\to0$.
Dla dostatecznie dużych $n$ mamy $|z_{n,k}|\leq1/2$
dla wszystkich $k$. Na tym kole określamy
$L(z)=\sum_{j=1}^\infty(-1)^{j+1}z^j/j$.
Różniczkując szereg potęgowy wewnątrz koła jednostkowego,
otrzymujemy $L'(z)=\sum_{j=0}^\infty(-z)^j=1/(1+z)$.
Pochodna funkcji $\exp L(z)/(1+z)$ jest więc zerowa,
a jej wartość w zerze wynosi jeden. Stąd $\exp L(z)=1+z$.
Z szeregu otrzymujemy
$|L(z)-z|\leq |z|^2/(1-|z|)\leq2|z|^2$.
W konsekwencji
\[
\left|\sum_k L(z_{n,k})-\sum_k z_{n,k}\right|
\leq 2\sum_k|z_{n,k}|^2\longrightarrow0.
\]
Nie trzeba utożsamiać sumy lokalnych logarytmów
z logarytmem głównym iloczynu. Z niezależności w każdym wierszu i własności
funkcji charakterystycznej sumy (rozdział~\ref{ch:expectation}) mamy
\[
\varphi_{\sum_k Z_{n,k}}(t)
=\prod_k(1+z_{n,k})
=\exp\left(\sum_k L(z_{n,k})\right)
\longrightarrow e^{-t^2/2}.
\]
Jest to funkcja charakterystyczna $\Normal(0,1)$,
ciągła w zerze. Twierdzenie Lévy'ego
(rozdział~\ref{ch:convergence}) kończy dowód.
\end{proof}

\begin{twierdzenie}[Twierdzenie Lindeberga--Fellera: pełna równoważność]
\label{thm:rp1-11-lindeberg-feller}
Przy założeniach schematu trójkątnego załóżmy dodatkowo
warunek Fellera:
\[
\max_{k\leq k_n}\frac{\Var(Y_{n,k})}{s_n^2}
\longrightarrow0.
\]
Wtedy zbieżność według rozkładu
$s_n^{-1}\sum_{k=1}^{k_n}Y_{n,k}\xrightarrow{d}\Normal(0,1)$
jest równoważna warunkowi Lindeberga.
\end{twierdzenie}
\begin{proof}
Kierunek wystarczający udowodniono wcześniej, bez
dodatkowego założenia Fellera; co więcej, warunek
Fellera wynika już z warunku Lindeberga (krok 1
poprzedniego dowodu).

Wykażemy konieczność. Niech $Z_{n,k}=Y_{n,k}/s_n$,
a $v_{n,k}=\E Z_{n,k}^2$. Mamy $\sum_k v_{n,k}=1$
i $\max_k v_{n,k}\to0$.
Wystarczy rozważyć funkcje charakterystyczne w jednym
punkcie, $t=1$. Połóżmy $z_{n,k}=\E e^{iZ_{n,k}}-1$. Z oszacowania Taylora
z poprzedniego dowodu i z $\E Z_{n,k}=0$ wynika
$|z_{n,k}|\leq v_{n,k}/2$. Zatem
\[
\max_k|z_{n,k}|\to0,\qquad
\sum_k |z_{n,k}|^2\leq
\frac14\sum_k v_{n,k}^2
\leq\frac14\max_k v_{n,k}\to0.
\]
Dla dostatecznie dużych $n$ możemy użyć tej samej
lokalnej gałęzi logarytmu $L$ co poprzednio.
Ponieważ $\exp L(z)=1+z$, mamy
$\log|1+z|=\operatorname{Re} L(z)$.
Z niezależności w wierszach i oszacowania
$|L(z)-z|\leq2|z|^2$ dostajemy
\[
\log\left|\E e^{i\sum_k Z_{n,k}}\right|
=\sum_k\operatorname{Re} z_{n,k}+r_n,
\qquad |r_n|\leq 2\sum_k|z_{n,k}|^2\to0.
\]
Założona zbieżność według rozkładu do $\Normal(0,1)$
implikuje zbieżność funkcji charakterystycznych
(rozdział~\ref{ch:convergence}), zatem lewa strona
dąży do $\log e^{-1/2}=-1/2$. Wobec tego
$\sum_k\E(1-\cos Z_{n,k})\to1/2$ i, ponieważ
$\sum_k\E Z_{n,k}^2=1$,
\[
\sum_k\E\left[\frac{Z_{n,k}^2}{2}
 -(1-\cos Z_{n,k})\right]\longrightarrow0.
\]

Dla rzeczywistego $u$ połóżmy
$h(u)=u^2/2-(1-\cos u)$. Funkcja jest nieujemna
i ściśle dodatnia poza zerem: dla $u\ne0$
mamy $1-\cos u=2\sin^2(u/2)<u^2/2$,
gdyż $|\sin v|<|v|$ dla $v\ne0$: dla $v>0$
mamy $|\int_0^v\cos t\,\dd\lambda_1(t)|
\leq\int_0^v|\cos t|\,\dd\lambda_1(t)<v$,
a dla $v<0$ korzystamy z nieparzystości sinusa.
Ponadto $h(u)/u^2\to1/2$ dla $|u|\to\infty$,
bo $0\leq1-\cos u\leq2$.
Dla każdego $\varepsilon>0$ ciągłość ilorazu
na $\{|u|\geq\varepsilon\}$, jego ścisła dodatniość
oraz dodatnia granica w nieskończoności dają
stałą $c_\varepsilon>0$, taką że
$h(u)\geq c_\varepsilon u^2$ dla $|u|\geq\varepsilon$.
Z nieujemności $h$ na całej prostej wynika
\[
0\leq\sum_k\E\left[
 Z_{n,k}^2\chi_{\{|Z_{n,k}|>\varepsilon\}}\right]
\leq \frac1{c_\varepsilon}\sum_k\E h(Z_{n,k})
\longrightarrow0.
\]
Jest to dokładnie warunek Lindeberga po normalizacji.
\end{proof}

\begin{wniosek}[Warunek Lapunowa]
Przy założeniach schematu trójkątnego załóżmy,
że dla pewnego $\delta>0$ wszystkie momenty
$\E|Y_{n,k}|^{2+\delta}$ są skończone i
\[
\frac{1}{s_n^{2+\delta}}
\sum_{k=1}^{k_n}\E|Y_{n,k}|^{2+\delta}
\longrightarrow0.
\]
Wówczas $s_n^{-1}\sum_{k=1}^{k_n}Y_{n,k}$
zbiega według rozkładu do $\Normal(0,1)$.
\end{wniosek}
\begin{proof}
Na zdarzeniu $\{|Y_{n,k}|>\varepsilon s_n\}$
mamy $|Y_{n,k}|^2\leq
|Y_{n,k}|^{2+\delta}/(\varepsilon s_n)^\delta$.
Sumując wartości oczekiwane i dzieląc przez $s_n^2$,
otrzymujemy lewą stronę warunku Lindeberga
nie większą niż $\varepsilon^{-\delta}$ razy
wyrażenie z założenia. To wyrażenie dąży do zera,
więc stosujemy poprzednie CTG.
\end{proof}

\begin{przyklad}[Zmienne niezależne, ale niejednakowo rozłożone]
Niech $X_k$ będą niezależne, $\E X_k=0$,
$|X_k|\leq C$ p.n. dla pewnej stałej $C>0$ i
$s_n^2=\sum_{k=1}^n\Var(X_k)$.
Jeśli $s_n\to\infty$, to dla ustalonego
$\varepsilon>0$ od pewnego $n$ mamy
$|X_k|\leq C<\varepsilon s_n$ dla każdego
$k\leq n$. Warunek Lindeberga jest więc
spełniony automatycznie, a zatem
$\sum_{k=1}^nX_k/s_n\xrightarrow d \Normal(0,1)$.
Wspólny rozkład nie jest potrzebny; ważna jest
całkowita wariancja oraz brak dominujących składników.
\end{przyklad}

\begin{przyklad}[Dlaczego warunek na duże składniki jest potrzebny]
W każdym wierszu niech $Y_{n,1}$ przyjmuje
wartości $\pm1$ z prawdopodobieństwami $1/2$,
a $Y_{n,k}=0$ dla $2\leq k\leq n$.
Wiersz jest niezależny, ma $s_n^2=1$,
lecz suma znormalizowana zawsze ma rozkład
Rademachera, a nie normalny.
Warunek Lindeberga nie zachodzi dla
$\varepsilon\in(0,1)$: jedyny niezerowy składnik
wnosi pełną jednostkę do sumy obciętych drugich
momentów.
\end{przyklad}

\begin{uwaga}[Dokładny zakres twierdzenia Lindeberga]
W twierdzeniu~\ref{thm:rp1-11-lindeberg-feller}
założenie Fellera jest istotne dla kierunku koniecznego.
Sama zbieżność sum do rozkładu normalnego nie implikuje
warunku Lindeberga. Istotnie, jeśli w każdym wierszu
$Y_{n,1}\sim\Normal(0,1)$, a wszystkie pozostałe
składniki są zerowe, to $s_n=1$, a suma ma dokładnie
rozkład $\Normal(0,1)$. Tymczasem dla każdego
$\varepsilon>0$ wyrażenie z warunku Lindeberga
jest równe dodatniej stałej
$\E[Y_{n,1}^2\chi_{\{|Y_{n,1}|>\varepsilon\}}]$,
niezależnej od $n$. W wersji twierdzenia o równoważności
potrzebny jest dodatkowy warunek Fellera:
$\max_{k\leq k_n}\Var(Y_{n,k})/s_n^2\to0$.
Warunek Lindeberga kontroluje nie tylko prawdopodobieństwo
rzadkiego dużego skoku, ale jego wkład do \emph{drugiego momentu}.
Zastąpienie warunku Lindeberga samym warunkiem
\[
\sum_{k=1}^{k_n}\Prob(|Y_{n,k}|>\varepsilon s_n)\to0
\qquad(\varepsilon>0)
\]
byłoby zbyt słabe. Wprawdzie warunek Lindeberga
implikuje powyższy warunek, gdyż
\[
\sum_k\Prob(|Y_{n,k}|>\varepsilon s_n)
\leq \frac1{\varepsilon^2 s_n^2}
\sum_k\E[Y_{n,k}^2\chi_{\{|Y_{n,k}|>\varepsilon s_n\}}],
\]
ale zdarzenie rzadkie może nieść niemal całą
wariancję sumy, jak pokazuje następny przykład.
\end{uwaga}

\begin{przyklad}[Rzadki skok o istotnej wariancji]
Dla każdego $n$ niech $Y_{n,1}$ przyjmuje $\sqrt n$
i $-\sqrt n$ z prawdopodobieństwami $1/(2n)$
każde, a zero z prawdopodobieństwem $1-1/n$.
Pozostałe wyrazy wiersza są zerowe. Mamy
$\E Y_{n,1}=0$, $s_n^2=\E Y_{n,1}^2=1$.
Realizujemy wszystkie wiersze na wspólnej przestrzeni
probabilistycznej (na przykład jako współrzędne przestrzeni
produktowej). Suma znormalizowana zbiega do zera według prawdopodobieństwa,
bo prawdopodobieństwo niezerowego wyniku wynosi $1/n$,
choć jej wariancja pozostaje równa jeden.
Dla dowolnego $\varepsilon>0$ i $n>\varepsilon^2$
wyrażenie z warunku Lindeberga wynosi $1$.
Jednocześnie suma prawdopodobieństw dużych skoków
wynosi $1/n$ i dąży do zera, więc słabszy warunek
z poprzedniej uwagi jest spełniony.
Nie można więc uzyskać granicy normalnej przez
argument opierający się wyłącznie na malejącej
częstości dużych skoków.

Nawet dołączenie warunku Fellera nie wystarcza do
naprawienia słabszego kryterium. Istotnie, w wierszu
o długości $n$ weźmy niezależne zmienne, z których
każda przyjmuje wartości $\pm\sqrt n$ z prawdopodobieństwem
$1/(2n^2)$ każda, a zero z prawdopodobieństwem $1-1/n^2$.
Każda ma średnią zero i wariancję $1/n$, więc
$s_n^2=1$ i warunek Fellera zachodzi. Dla ustalonego
$\varepsilon>0$ i $n>\varepsilon^2$ suma
prawdopodobieństw dużych skoków wynosi $1/n$,
natomiast suma obciętych drugich momentów wynosi $1$.
Ponadto z nierówności dla prawdopodobieństwa sumy
zdarzeń otrzymujemy
\[
\Prob\left(\left|\sum_{k=1}^n Y_{n,k}\right|>\varepsilon\right)
\leq\Prob\left(\bigcup_{k=1}^n\{Y_{n,k}\ne0\}\right)
\leq\frac1n.
\]
Suma zbiega więc według prawdopodobieństwa do zera,
a nie według rozkładu do $\Normal(0,1)$.
\end{przyklad}

\section{Metoda delty: granice nieliniowych funkcji średnich}

Centralne twierdzenie graniczne mówi o przybliżeniu
normalnym dla średniej. Jeśli interesuje nas
$f(\overline X_n)$, np. logarytm, odwrotność albo
pierwiastek, nie powinniśmy od nowa analizować
całego rozkładu. Wystarczy lokalne rozwinięcie
funkcji wokół granicy średnich.

\begin{twierdzenie}[Metoda delty pierwszego rzędu]
Niech $T_n$ będą rzeczywistymi zmiennymi losowymi,
$\theta\in\R$, $a_n>0$ i $a_n\to\infty$.
Załóżmy, że $a_n(T_n-\theta)\xrightarrow d Z$
dla pewnej rzeczywistej zmiennej losowej $Z$.
Jeśli $g:\R\to\R$ jest funkcją borelowską, różniczkowalną
w punkcie $\theta$, to
\[
a_n\bigl(g(T_n)-g(\theta)\bigr)
\xrightarrow d g'(\theta)Z.
\]
\end{twierdzenie}
\begin{proof}
Najpierw zauważmy, że rodzina rozkładów
$a_n(T_n-\theta)$ jest ciasna, bo zbiega według
rozkładu (rozdział~\ref{ch:convergence}). Dla dowolnych $\varepsilon,\delta>0$ z ciasności
wybieramy $M>0$ tak, by
$\sup_n\Prob(|a_n(T_n-\theta)|>M)<\delta$.
Dla dostatecznie dużych $n$ zachodzi
$\varepsilon a_n\geq M$, zatem
\[
\Prob(|T_n-\theta|>\varepsilon)
=\Prob(|a_n(T_n-\theta)|>\varepsilon a_n)
\leq\Prob(|a_n(T_n-\theta)|>M)<\delta.
\]
Dowolność $\delta$ dowodzi, że $T_n-\theta$
zbiega do zera według prawdopodobieństwa,
czyli $T_n\xrightarrow{\Prob}\theta$.
Z różniczkowalności $g$ w $\theta$ wynika
reprezentacja
$g(x)-g(\theta)=(x-\theta)(g'(\theta)+r(x))$,
gdzie $r(x)\to0$ dla $x\to\theta$
i przyjmujemy $r(\theta)=0$.
Funkcja $r$ jest borelowska, a z $r(x)\to0$ przy
$x\to\theta$ i $r(\theta)=0$ wynika, że dla każdego
$\eta>0$ istnieje $\rho>0$ takie, że
$\{|r(T_n)|>\eta\}\subseteq\{|T_n-\theta|\geq\rho\}$.
Stąd $r(T_n)\xrightarrow{\Prob}0$.
Po pomnożeniu przez $a_n$ mamy
\[
a_n(g(T_n)-g(\theta))
=[a_n(T_n-\theta)]\,[g'(\theta)+r(T_n)].
\]
Twierdzenie Slutsky'ego (rozdział~\ref{ch:convergence})
daje granicę $Z\,g'(\theta)$, zgodnie z tezą.
\end{proof}

\begin{wniosek}[Metoda delty drugiego rzędu]
Przy założeniach poprzedniego twierdzenia załóżmy
dodatkowo, że $g$ jest różniczkowalna w otoczeniu $\theta$,
jej pochodna jest różniczkowalna w $\theta$ oraz
$g'(\theta)=0$.
Wtedy
\[
a_n^2(g(T_n)-g(\theta))
\xrightarrow d\frac{g''(\theta)}2Z^2.
\]
\end{wniosek}
\begin{proof}
Uzasadnimy rozwinięcie Taylora z resztą Peana
przy podanej, lokalnej hipotezie. Niech
$h(x)=g(x)-g(\theta)-\frac12g''(\theta)(x-\theta)^2$.
Ponieważ $g'(\theta)=0$ i $g'$ jest różniczkowalna
w $\theta$, mamy $h'(x)=o(|x-\theta|)$ dla $x\to\theta$.
Dla każdego $\eta>0$ w dostatecznie małym otoczeniu
$\theta$ zachodzi więc $|h'(x)|\leq\eta|x-\theta|$.
Z twierdzenia Lagrange'a o wartości średniej,
zastosowanego do $h$ między $\theta$ i $x$, wynika
$|h(x)|\leq\eta|x-\theta|^2$, ponieważ $h(\theta)=0$.
Stąd
$g(x)-g(\theta)=\frac12g''(\theta)(x-\theta)^2
+o((x-\theta)^2)$. Ponieważ $T_n\to\theta$
według prawdopodobieństwa, iloraz reszty przez
$(T_n-\theta)^2$ (z wartością zero na zdarzeniu
$\{T_n=\theta\}$) dąży do zera według prawdopodobieństwa
na mocy tego samego argumentu co w dowodzie
metody delty pierwszego rzędu.
Zatem
$a_n^2(g(T_n)-g(\theta))
=(\frac12g''(\theta)+o_{\Prob}(1))
[a_n(T_n-\theta)]^2$,
gdzie $o_{\Prob}(1)$ oznacza zmienną losową zależną od $n$,
która zbiega do zera według prawdopodobieństwa.
Z twierdzenia o ciągłym przekształceniu
$[a_n(T_n-\theta)]^2\xrightarrow d Z^2$;
Slutsky kończy dowód.
\end{proof}

\begin{przyklad}[Asymptotyka logarytmu średniej]
Niech $X_1,X_2,\ldots$ będą niezależne, jednakowo
rozłożone, $\E X_1=m>0$ i
$0<\Var(X_1)=\sigma^2<\infty$.
Z prawa wielkich liczb $\overline X_n\to m$
według prawdopodobieństwa. Ponieważ
$\{\overline X_n\leq0\}\subseteq
\{|\overline X_n-m|\geq m\}$ i $m>0$,
prawdopodobieństwo zdarzenia $\{\overline X_n>0\}$
dąży do jedności;
definiujemy $T_n=\log\overline X_n$ na zdarzeniu
$\{\overline X_n>0\}$, a $T_n=\log m$ na jego dopełnieniu.
Z klasycznego CTG mamy
$\sqrt n(\overline X_n-m)\xrightarrow d\Normal(0,\sigma^2)$.
Z metody delty dla $g(x)=\log x$ w otoczeniu $m$
(z borelowskim rozszerzeniem
$g(x)=\log m$ dla $x\leq0$) i z $g'(m)=1/m$
otrzymujemy
$\sqrt n(T_n-\log m)\xrightarrow d
\Normal(0,\sigma^2/m^2)$.
Wzór pokazuje, skąd bierze się czynnik $1/m$
w przybliżeniu niepewności logarytmu średniej.
\end{przyklad}

\begin{przyklad}[Gdy pochodna pierwszego rzędu znika]
Niech $Z_1,Z_2,\ldots$ będą niezależnymi zmiennymi
losowymi o rozkładzie $\Normal(0,1)$ i niech
$T_n=\overline Z_n$.
Niech ponadto $Z\sim\Normal(0,1)$ oznacza zmienną
o rozkładzie granicznym. Mamy $\sqrt n T_n\xrightarrow d \Normal(0,1)$.
Dla $g(x)=x^2$ zachodzi $g'(0)=0$
i $g''(0)=2$, zatem
$nT_n^2\xrightarrow d Z^2\sim\chi^2_1$,
gdzie rozkład kwadratu standardowej zmiennej normalnej
wyznaczono w rozdziale~\ref{ch:distributions}.
W tym przykładzie równość rozkładów jest nawet
prawdziwa dla każdego $n$, ponieważ
$\sqrt n\,T_n\sim \Normal(0,1)$: niezależność
daje funkcję charakterystyczną
$\prod_{k=1}^n e^{-t^2/(2n)}=e^{-t^2/2}$,
a jednoznaczność funkcji charakterystycznej
identyfikuje rozkład (rozdział~\ref{ch:convergence}).
Zastosowanie
metody delty pierwszego rzędu dawałoby tylko
zdegenerowaną granicę zero przy skali $\sqrt n$.
\end{przyklad}

\section{Studentyzacja: nieznana wariancja w CTG}

\begin{twierdzenie}[CTG po zastąpieniu wariancji estymatorem]
Niech $X_1,X_2,\ldots$ będą niezależnymi,
jednakowo rozłożonymi rzeczywistymi zmiennymi losowymi,
$\E X_1=m$ i $0<\Var(X_1)=\sigma^2<\infty$.
Dla $n\geq2$ połóżmy
$V_n=n^{-1}\sum_{k=1}^n(X_k-\overline X_n)^2$.
Wtedy $V_n\to\sigma^2$ prawie na pewno i
\[
\frac{\sqrt n(\overline X_n-m)}{\sqrt{V_n}}
\xrightarrow d\Normal(0,1),
\]
przy dowolnym mierzalnym, rzeczywistym określeniu ilorazu na zdarzeniu
$\{V_n=0\}$.
\end{twierdzenie}
\begin{proof}
Rozwijając kwadrat, otrzymujemy dokładną tożsamość
$V_n=n^{-1}\sum_{k\leq n}X_k^2-\overline X_n^2$.
Ponieważ $\E X_1^2<\infty$, mocne prawo wielkich
liczb można zastosować zarówno do $X_k$, jak i $X_k^2$:
$\overline X_n\to m$ i
$n^{-1}\sum_{k\leq n}X_k^2\to\E X_1^2$
prawie na pewno. Stąd $V_n\to\sigma^2>0$ p.n.,
a więc na zbiorze pełnej miary od pewnego indeksu
$V_n>0$. Definiując $R_n=1/\sqrt{V_n}$ dla $V_n>0$
i $R_n=0$ poza tym zdarzeniem, mamy
$R_n\to1/\sigma$ prawie na pewno.
Ponadto $\Prob(V_n=0)\to0$ ze zbieżności według
prawdopodobieństwa do dodatniej stałej $\sigma^2$.
Z podstawowego CTG (rozdział~\ref{ch:clt})
$\sqrt n(\overline X_n-m)\xrightarrow d \Normal(0,\sigma^2)$.
Twierdzenie Slutsky'ego daje wskazaną granicę dla
iloczynu $\sqrt n(\overline X_n-m)R_n$. Każde inne
mierzalne, rzeczywiste określenie ilorazu na
$\{V_n=0\}$ daje zmienną $W_n$, której różnica
z iloczynem $A_n=\sqrt n(\overline X_n-m)R_n$
spełnia dla każdego $\varepsilon>0$
\[
\Prob(|W_n-A_n|>\varepsilon)\leq\Prob(V_n=0)\to0.
\]
Ponowne zastosowanie twierdzenia Slutsky'ego do
$W_n=A_n+(W_n-A_n)$ daje tę samą granicę według rozkładu.
\end{proof}

\begin{uwaga}
\emph{Asymptotyczna} normalność zmiennej
studentyzowanej zachodzi dla wielu rozkładów
z drugim momentem. Nie jest tym samym co
\emph{dokładny} rozkład Studenta, który
wynika z dodatkowych własności próby normalnej.
Przy zwykłej wariancji z próby
$S_n^2=(n-1)^{-1}\sum_{k=1}^n(X_k-\overline X_n)^2$
mamy $S_n^2=nV_n/(n-1)$, więc zastąpienie $V_n$
przez $S_n^2$ nie zmienia granicy asymptotycznej.
Dokładny rozkład Studenta o $n-1$ stopniach swobody
dla próby normalnej dotyczy ilorazu z $S_n$, nie
z $\sqrt{V_n}$. Uzasadnimy to dla każdego $n\geq2$.

Jeśli $X_k\sim\Normal(m,\sigma^2)$, połóżmy
$Z_k=(X_k-m)/\sigma$ i niech $Z=(Z_1,\ldots,Z_n)$.
Niezależność daje gęstość łączną
$(2\pi)^{-n/2}\exp(-\|z\|^2/2)$.
Wybierzmy macierz ortogonalną $O$ o pierwszym wierszu
$n^{-1/2}(1,\ldots,1)$ (uzupełniamy ten wektor do bazy
ortonormalnej). Dla $W=OZ$ wzór na gęstość
przekształcenia liniowego daje tę samą gęstość,
ponieważ $|\det O|=1$ i $O$ zachowuje normę.
Gęstość jest iloczynem gęstości standardowych
rozkładów normalnych, zatem $W_1,\ldots,W_n$ są
niezależne i mają rozkład $\Normal(0,1)$.
Ponadto
\[
W_1=\frac{\sqrt n(\overline X_n-m)}{\sigma},
\qquad
\sum_{j=2}^n W_j^2
=\sum_{k=1}^n Z_k^2-W_1^2
=\frac{(n-1)S_n^2}{\sigma^2}.
\]
Ostatnia równość wynika z rozwinięcia
$\sum_k(X_k-\overline X_n)^2
=\sum_k(X_k-m)^2-n(\overline X_n-m)^2$.
Zatem $U=\sum_{j=2}^n W_j^2\sim\chi^2_{n-1}$
i $U$ jest niezależna od $W_1$. Mamy $U>0$ p.n.,
gdyż $\Prob(W_2=0)=0$. Z definicji rozkładu Studenta
(rozdział~\ref{ch:distributions}) otrzymujemy
\[
\frac{\sqrt n(\overline X_n-m)}{S_n}
=\frac{W_1}{\sqrt{U/(n-1)}}\sim t_{n-1}.
\]
Iloraz z $\sqrt{V_n}$ ma rozkład tego ilorazu
pomnożonego przez $\sqrt{n/(n-1)}$.
Rozróżnienie to jest ważne przy przejściu do
statystyki matematycznej w tomie IV.
\end{uwaga}

\section{Jednostajna zbieżność dystrybuanty empirycznej}

Dla próby $X_1,\ldots,X_n$ obserwujemy nie tylko
średnią, ale i empiryczny rozkład: każda obserwacja
otrzymuje masę $1/n$. Dystrybuanta empiryczna ma
postać
$F_n(x)=n^{-1}\sum_{k=1}^n\chi_{\{X_k\leq x\}}$.
Dla ustalonego $x$ jest to średnia prób Bernoulliego
z parametrem $F(x)$, więc MPWL daje
$F_n(x)\to F(x)$ p.n., o ile obserwacje są niezależne
i jednakowo rozłożone ze wspólną dystrybuantą $F$.
Istotne jest mocniejsze pytanie:
czy możemy jednocześnie kontrolować błąd
na \emph{całej prostej}? W tym miejscu potrzebujemy
odrobiny więcej niż punktowego MPWL.

\begin{twierdzenie}[Glivenki--Cantellego]
Niech $X_1,X_2,\ldots$ będą niezależnymi
jednakowo rozłożonymi rzeczywistymi zmiennymi losowymi
o wspólnej dystrybuancie $F$ i niech
$F_n(x)=n^{-1}\sum_{k=1}^n\chi_{\{X_k\leq x\}}$.
Wówczas
\[
\sup_{x\in\R}|F_n(x)-F(x)|\longrightarrow0
\qquad\text{prawie na pewno}.
\]
\end{twierdzenie}
\begin{proof}
\emph{Krok pierwszy: rozkład jednostajny.}
Na pomocniczej przestrzeni $(0,1)^{\N}$ z produktowym
$\sigma$-ciałem borelowskim i miarą produktową
$\bigotimes_{k\geq1}\lambda_1|_{(0,1)}$
wybieramy odwzorowania współrzędnych $U_k$.
Z własności miary produktowej są one niezależne
i mają rozkład $\Unif(0,1)$; w każdym punkcie
przestrzeni przyjmują wartości w $(0,1)$. Niech
$H_n(t)=n^{-1}\sum_{k\leq n}\chi_{\{U_k\leq t\}}$.
Dla każdego ustalonego $t\in[0,1]$ zmienne
$\chi_{\{U_k\leq t\}}$ są niezależne i jednakowo
rozłożone, o średniej $t$. MPWL
(rozdział~\ref{ch:lln}) daje $H_n(t)\to t$ p.n. Zastosujemy je
jednocześnie tylko do przeliczalnego zbioru
liczb wymiernych $t\in[0,1]$. Ich przecięcie
zdarzeń pełnej miary nadal ma miarę jeden.
Ustalmy trajektorię z tego przecięcia i
liczbę całkowitą $m\geq1$. Dla
$0\leq j\leq m$ mamy
$H_n(j/m)\to j/m$.
Dla każdego $t\in[0,1]$ wybieramy
$j\in\{0,\ldots,m-1\}$ z $j/m\leq t\leq(j+1)/m$.
Monotoniczność
dystrybuanty empirycznej daje
\[
H_n(j/m)-\frac{j+1}{m}\leq H_n(t)-t
\leq H_n((j+1)/m)-\frac{j}{m}.
\]
Zatem
$\sup_{t\in[0,1]}|H_n(t)-t|\leq
\max_{0\leq j\leq m}|H_n(j/m)-j/m|+1/m$.
Przechodząc do granicy górnej przy $n\to\infty$,
otrzymujemy oszacowanie przez $1/m$ dla każdego $m$,
gdyż maksimum skończonej liczby błędów w punktach
siatki dąży do zera.
Ponieważ $m$ jest dowolne, wynika stąd
$\sup_{t\in[0,1]}|H_n(t)-t|\to0$ p.n.

\emph{Krok drugi: dowolna dystrybuanta.}
Niech $Q(u)=\inf\{x\in\R:F(x)\geq u\}$,
$u\in(0,1)$, będzie uogólnioną funkcją odwrotną
dystrybuanty (funkcją kwantylową).
Ponieważ $F(-\infty)=0$ i $F(+\infty)=1$, zbiór
$\{x:F(x)\geq u\}$ jest niepusty i ograniczony z dołu,
więc $Q(u)\in\R$.
Z definicji kresu dolnego wybieramy punkty tego
zbioru z przedziału $[Q(u),Q(u)+1/j)$ dla $j\geq1$.
Prawostronna ciągłość i monotoniczność $F$ dają
$F(Q(u))\geq u$. Stąd
\[
Q(u)\leq x\quad\Longleftrightarrow\quad u\leq F(x).
\]
Implikacja w prawo wynika z monotoniczności $F$,
a w lewo z definicji kresu dolnego. Funkcja $Q$
jest niemalejąca, a powyższa równoważność
pokazuje również jej borelowskość:
$\{u:Q(u)\leq x\}=(0,1)\cap(0,F(x)]$.
Argument nie wymaga ciągłości $F$ z lewej strony
ani jej ścisłej monotoniczności.
Wobec tego zmienne $Q(U_k)$ są niezależne jako
mierzalne funkcje niezależnych zmiennych oraz
$\Prob(Q(U_k)\leq x)=\Prob(U_k\leq F(x))=F(x)$.
Dla skonstruowanej próby oznaczmy dystrybuantę
empiryczną przez $\widetilde F_n$. Równoważność
dotycząca $Q$ zachodzi punktowo dla wszystkich $x$,
więc $\widetilde F_n(x)=H_n(F(x))$.
Dlatego
$\sup_x|\widetilde F_n(x)-F(x)|
\leq\sup_{t\in[0,1]}|H_n(t)-t|\to0$
prawie na pewno w skonstruowanym modelu.

Ostatni krok przenosimy do dowolnej niezależnej próby
o tej samej dystrybuancie. Oba ciągi mają ten sam
rozkład na $\R^{\N}$ wyposażonej w produktowe
$\sigma$-ciało borelowskie: na zbiorach cylindrycznych
zależnych od skończenie wielu współrzędnych ich
rozkłady są równe z niezależności, a zbiory te tworzą
$\pi$-układ generujący to $\sigma$-ciało. Jednoznaczność
miary probabilistycznej na takim układzie wynika
z twierdzenia o $\pi$- i $\lambda$-układach
(por. argument o jednoznaczności miary na cylindrach
w rozdziale~\ref{ch:probability}).
Dla ciągu $z=(z_k)_{k\geq1}$ połóżmy
\[
D_n(z)=\sup_{q\in\mathbb Q}
\left|\frac1n\sum_{k=1}^n\chi_{(-\infty,q]}(z_k)
-F(q)\right|.
\]
Jest to funkcja mierzalna jako przeliczalne supremum
funkcji mierzalnych. Z prawostronnej ciągłości obu
dystrybuant supremum po liczbach wymiernych jest
równe supremum po wszystkich liczbach rzeczywistych:
dla dowolnego $x$ wybieramy $q_j\downarrow x$.
Zdarzenie
\[
\{z:D_n(z)\to0\}
=\bigcap_{m\geq1}\bigcup_{N\geq1}
\bigcap_{n\geq N}\{z:D_n(z)<1/m\}
\]
należy więc do produktowego $\sigma$-ciała.
Ma ono jednakowe prawdopodobieństwo w obu modelach,
a w modelu kwantylowym prawdopodobieństwo to jest
równe jeden. Dowodzi to tezy na pierwotnej przestrzeni
probabilistycznej.
\end{proof}

Rysunek~\ref{fig:rp1-11-empirical-cdf} ilustruje
jednostajną odległość dystrybuanty empirycznej
od dystrybuanty teoretycznej na dwóch skończonych próbach.

\begin{figure}[htbp]
\centering
\begin{tikzpicture}[x=4.6cm,y=4.6cm,>=Latex]
\foreach \shift/\lab in {0cm/5,6.808cm/20}{
  \begin{scope}[xshift=\shift]
    \draw[->] (0,0)--(1.09,0) node[right] {$t$};
    \draw[->] (0,0)--(0,1.09) node[above] {$H_n(t)$};
    \foreach \x/\tick in {.5/{0{,}5},1/1}
      {\draw[gray!20] (\x,0)--(\x,1);
       \draw (\x,0)--++(0,-.018) node[below] {$\tick$};}
    \foreach \y/\tick in {.5/{0{,}5},1/1}
      {\draw[gray!20] (0,\y)--(1,\y);
       \draw (0,\y)--++(-.018,0) node[left] {$\tick$};}
    \draw[gray!85,dashed,thick] (0,0)--(1,1);
    \node[above] at (.5,1.13) {$n=\lab$};
  \end{scope}
}
\draw[blue!75!black,very thick] (0,0)--(0.12,0)--(0.12,0.2)--(0.34,0.2)--(0.34,0.4)--(0.47,0.4)--(0.47,0.6)--(0.78,0.6)--(0.78,0.8)--(0.91,0.8)--(0.91,1)--(1,1);
\begin{scope}[xshift=6.808cm]
\draw[blue!75!black,very thick] (0,0)--(0.03,0)--(0.03,0.05)--(0.08,0.05)--(0.08,0.1)--(0.12,0.1)--(0.12,0.15)--(0.18,0.15)--(0.18,0.2)--(0.24,0.2)--(0.24,0.25)--(0.29,0.25)--(0.29,0.3)--(0.34,0.3)--(0.34,0.35)--(0.39,0.35)--(0.39,0.4)--(0.43,0.4)--(0.43,0.45)--(0.47,0.45)--(0.47,0.5)--(0.51,0.5)--(0.51,0.55)--(0.55,0.55)--(0.55,0.6)--(0.62,0.6)--(0.62,0.65)--(0.66,0.65)--(0.66,0.7)--(0.72,0.7)--(0.72,0.75)--(0.78,0.75)--(0.78,0.8)--(0.85,0.8)--(0.85,0.85)--(0.91,0.85)--(0.91,0.9)--(0.95,0.9)--(0.95,0.95)--(0.98,0.95)--(0.98,1)--(1,1);
\end{scope}
\draw[blue!75!black,very thick] (.15,-.20)--(.27,-.20)
node[right,black] {empiryczna};
\draw[gray!85,dashed,thick] (1.48,-.20)--(1.60,-.20)
node[right,black] {teoretyczna $F(t)=t$};
\end{tikzpicture}
\caption{Dystrybuanty empiryczne rozkładu jednostajnego
na $(0,1)$ dla dwóch zagnieżdżonych prób ilustracyjnych:
pierwszych pięciu oraz wszystkich dwudziestu obserwacji.
Błąd jednostajny w tym przykładzie wynosi odpowiednio
$0{,}18$ oraz $0{,}06$.
Pionowe odcinki są konwencją graficzną oznaczającą skoki;
wartość w punkcie skoku jest wartością prawostronną.
Twierdzenie Glivenki--Cantellego nie orzeka, że błąd
maleje monotonicznie dla każdej realizacji.}
\label{fig:rp1-11-empirical-cdf}
\end{figure}

\begin{przyklad}[Dystrybuanta empiryczna próby Bernoulliego]
Niech $X_k\sim\operatorname{Bernoulli}(p)$ będą
niezależne i przyjmują wartości w $\{0,1\}$ (z dokładnością
do wspólnego zbioru zerowego). Dla $x<0$ obie dystrybuanty są zerowe,
a dla $x\geq1$ obie są jedynką. Na $0\leq x<1$
mamy $F_n(x)=1-\overline X_n$ oraz $F(x)=1-p$.
Stąd
$\sup_x|F_n(x)-F(x)|=|\overline X_n-p|$.
W tym szczególnym przypadku twierdzenie
Glivenki--Cantellego jest dokładnie mocnym prawem
wielkich liczb dla częstości sukcesów.
\end{przyklad}

\section{Koncentracja i funkcje Lipschitza}

\begin{przyklad}[Koncentracja odpowiedzi układu, a nie tylko średniej]
Niech $B_1,\ldots,B_n$ będą niezależnymi próbami Bernoulliego
z tym samym parametrem $p\in[0,1]$,
$S_n=\sum_{k=1}^n B_k$, a $g:\R\to\R$ będzie
funkcją $L$-Lipschitza z $L>0$, tzn.
$|g(x)-g(y)|\leq L|x-y|$ dla wszystkich $x,y$.
Z udowodnionej wcześniej nierówności Hoeffdinga
wynika dla każdego $\varepsilon>0$
\[
\Prob\bigl(|g(S_n/n)-g(p)|\geq L\varepsilon\bigr)
\leq 2e^{-2n\varepsilon^2}.
\]
Ograniczenie $L>0$ jest konieczne przy zapisanym
zdarzeniu z nierównością $\geq$: dla $L=0$
zdarzenie $\{|g(S_n/n)-g(p)|\geq0\}$ jest pewne,
a prawa strona może być mniejsza od jedności.
Dla $L>0$ wynik wynika z implikacji
$|g(S_n/n)-g(p)|\geq L\varepsilon\Longrightarrow
|S_n/n-p|\geq\varepsilon$.
Ogólniejsze narzędzia koncentracji dla funkcji wielu
niezależnych zmiennych, w tym nierówność ograniczonych
różnic McDiarmida opartą na nierówności
Azumy--Hoeffdinga, przedstawimy w rozdziale~\ref{ch:martyngaly}.
\end{przyklad}

\section{Synteza: jak wybrać właściwe twierdzenie?}

Jeśli interesuje nas wartość średniej przy dużym $n$,
stosujemy SPWL lub MPWL. Jeśli potrzebujemy rozkładu
fluktuacji średniej o skali $n^{-1/2}$ przy niezależnych,
jednakowo rozłożonych składnikach o dodatniej
skończonej wariancji, używamy klasycznego CTG.
Dla niejednakowych składników stosujemy odpowiednią
normalizację przez $s_n$ oraz warunek Lindeberga lub
Lapunowa. Gdy pytamy o konkretną liczbę prób potrzebną
do ograniczenia ryzyka odchylenia, pomocna jest
nierówność Hoeffdinga. Metoda delty opisuje
przekształcone estymatory, studentyzacja pozwala
zastąpić nieznaną wariancję jej empiryczną wersją,
a twierdzenie Glivenki--Cantellego kontroluje
cały rozkład empiryczny, nie tylko jedną średnią.
W przypadku niejednakowych składników warunek Fellera
wraz z granicą normalną jest równoważny warunkowi
Lindeberga; sam zanik prawdopodobieństw dużych skoków
nie zapewnia takiej granicy.
Wszystkie te rezultaty będą miały dalsze zastosowania
w teorii procesów stochastycznych i statystyce.

